\thispagestyle{empty}
 \begin{center}
 Міністерство освіти і науки України\\
 Національний університет «Києво-Могилянська 
академія»\\
Факультет інформатики\\
Кафедра математики
\end{center}
\begin{flushright}
ЗАТВЕРДЖУЮ\\
Зав. кафедри математики,\\
доцент, кандидат фіз.-мат. наук\\ 
$\underset{\text{(підпис)}}{\underline{\hspace{3.6cm}}}$\emph{Чорней Р.К.}\\
\begin{tabular}{l}
 ``\noindent\rule{1cm}{0.4pt}''\noindent\rule{3.9cm}{0.4pt} 2025\\
\end{tabular}
\end{flushright}
\vspace{1mm}
\begin{center}
ІНДИВІДУАЛЬНЕ ЗАВДАННЯ\\
на кваліфікаційну роботу\\
студенту 4-го курсу факультету інформатики\\
Нанарову Іллі Олександровичу
\end{center}
\noindent{\bf Тема:} «Дослідження якості україно-англійського синхронного перекладу мовлення: семантична точність та міжфразова когерентність»\\
\noindent{\bf Зміст кваліфікаційної роботи:}\\
\hspace*{13mm}
\begin{tabular}{l}
 Анотація\\
 1. Вступ\\
 2. Теоретичні основи: моделі ASR, машинного перекладу,\\
 методи оцінювання якості\\
 3. Корпус дослідження: формування та характеристики матеріалу\\
 4. Архітектура та компоненти системи: ASR (Whisper), MT\\
 (OPUS-MT, GPT-4o, Claude, Gemini), система метрик, дизайн експерименту\\
 5. Результати: оцінка якості ASR та MT, каскадний ефект,\\
 порівняння з LLM-системами (GPT-4o, Claude, Gemini)\\
 Висновки\\
 Список літератури \end{tabular}
\vspace{5mm} 
\begin{flushright}
 
 \begin{tabular}{clclc} 
  \noindent Дата видачі ``\noindent\rule{1cm}{0.4pt}''\noindent\rule{2.5cm}{0.4pt} 2025 & Керівник  $\underset{\text{(підпис)}}{\underline{\hspace{3.6cm}}}$\\
  
  & Завдання отримав  $\underset{\text{(підпис)}}{\underline{\hspace{3.6cm}}}$\\
 \end{tabular}
\end{flushright}